\documentclass[11pt,fleqn]{report}

\usepackage{cml}
\usepackage{amssymb}
\usepackage{graphicx}

\newcommand{\ModelName}{Polygon Enclosure}
\newcommand{\ModelAuthor}{Gary Turner}
\newcommand{\ModelAffiliation}{Odyssey Space Research}
\newcommand{\ModelDate}{May 2024}

\setcounter{tocdepth}{3}
\setcounter{secnumdepth}{3}

\begin{document}
\pagestyle{empty}
\titlePage
\pagenumbering{roman}
\tableofcontents % Print table of contents
\vfill
{\Large \textbf {Revision History}}
{\large
\begin{center}
\begin{tabular}{|l||l|l|l|}
\hline
 \textbf{Version} & \textbf{Date} & \textbf{Author} & \textbf{Purpose} \\
\hline \hline
1 & May 2024 & Gary Turner & Initial Version \\ \hline
\end{tabular}
\end{center}
}


\chapter{Introduction} % Make a "chapter" heading
\pagestyle{plain}
\pagenumbering{arabic} % Start text with arabic 1
The Polygon Enclosure model provides a mechanism for determining whether a point
lies within a bounding polygon.

There are two implementations: one for a flat Euclidean plane and one for a
spherical surface. At the time of writing, each implementation had slightly
different target requirements and the implementations were tailored to those
specific purposes. Consequently, there are substantive differences between the
two implementations that go beyond the difference in geometry.
At some time, it may be desirable to implement both geometries with a consistent
underlying architecture, but at the time this model was written, there was not
time or sufficient priority to do so. This is left as a future exercise.



\chapter{Requirements}
\begin{enumerate}
  \item The model shall quickly and efficiently identify whether a point on a
    Euclidean plane lies within a convex polygon prescribed within that plane.
  \begin{itemize}
    \item The model shall provide an input interface to support user-definition
      of the posiitons of the polygon vertices.
    \item The algorithm may assume the polygon is convex, but must verify the
      user-defined vertices form a convex polygon if such an assumption is
      used.
  \end{itemize}
  \item The model shall quickly and efficiently identify whether each of a
    set of points on the surface of a unit-sphere lie within a convex
    polygon prescribed on that spherical surface.
  \begin{itemize}
    \item The model shall provide an input interface to support user-definition
      of the posiitons of the polygon vertices in both Cartesian
      coordinates and angular coordinates.
    \item The algorithm may assume the polygon is convex, but must verify the
      user-defined vertices form a convex polygon if such an assumption is
      used.
  \end{itemize}
\end{enumerate}

%******************************************************************************
%******************************************************************************
%******************************************************************************
%******************************************************************************
\chapter{Model Specification}
\section{Code Structure}
\subsection{Polygon Enclosure (Planar)}
Signature: \texttt{PolygonEnclosure::in\_polygon()}

The planar implementation is provided by the class-template
\textit{PolygonEnclosure}.

\subsection{Polygon Enclosure (Unit Sphere)}
Signature: \texttt{PolygonEnclosureSphere::in\_polygon( double position[3])}

The spherical implementation is provided by the class-template
\textit{PolygonEnclosureSphere}.

%******************************************************************************
%******************************************************************************
%******************************************************************************
\section{Mathematical Formulation}

\subsection{Polygon Enclosure (Planar)}

The implementation works with an operational assumption that the vertices of
the polygon are defined and fixed within the plane, and we routinely wish to
identify whether a dynamic point of interest lies within that fixed polygon.
Because we are working on a plane, we can consider
our points -- for the purposes of this derivation -- as lying in the x-y
plane.

The vertices of the polygon are defined in a sequential order, with the
perimeter of the polygon being that sequence of lines joining adjacent
vertices. The line joining the first and last vertex completes the
perimeter.

At model initialization, the the $N$ vertices are tested to ensure that they
form a convex polygon (i.e. no interior angles are greater than 180
degrees).

To verify the convex nature, we take the cross-product between adjoining
line segments, and consider the sign of the z-component of the resulting
vector (the x- and y-components will be 0 by definition of cross product of
two co-planar vectors). If the sign of the z-component is not consistent
across all adjacent line-segment pairings, the polygon is not convex, and
cannot be used for this model.

\begin{figure}[h!]
  \begin{center}
    \label{fig:polygon_fault}
    \includegraphics[scale=0.5]{Figs/non_convex_polygon.png}
    \caption{Illustration of the detection of a non-convex polygon}
  \end{center}
\end{figure}

As an illustrative example, in the diagram below, $\vec a \times \vec b$
would be ``out of the page'', as would $\vec b \times \vec c$.
However, $\vec c \times \vec d$ would be ``into the page''. This sign
change means the polygon is not convex. Note that having three adjacent
vertices on a line does not violate the convex test, but it is both
redundant (the middle point can be eliminated) and risky (if the middle
point numerically rounds to the wrong side of that line, the entire polygon
will be rejected).

Also during initialization, we evaluate the minimum and maximum x- and y-
coordinate values of all vertices and can use these to draw a rectangle
around the polygon.
\begin{figure}[h!]
  \begin{center}
    \label{fig:polygon_convex}
    \includegraphics[scale=0.5]{Figs/polygon_convex_with_box.png}
    \caption{Illustration of a convex polygon and its bounding box.}
  \end{center}
\end{figure}

With the polygon initialized, we can now apply a similar algorithm as
necessary to identify whether a point lies within the polygon. First though,
if the point is outside the bounding box, it is outside the polygon. This is
an easy check to make and saves testing multiple vertices. If it is inside
the bounding box, we need to test each vertex.

Consider the point of interest at $\vec r$; we can generate a vector
$\vec r_i = \vec R_i - \vec r$ from vertex $i$ (at position $\vec R_i$)
to the point of interest for each vertex.

If the cross product of these vectors from adjacent vertices produces the
same sign as that identified during the convex test, then the point is
``inside'' the polygon edge joining those two vertices.
If the point is ``inside'' all polygon edges, it is inside the polygon.
Otherwise, it is outside the polygon.

Rather than computing the relative coordinates from each vertex to the point
of interest, we can simplify this process. Note that:
\begin{equation*}
  (\vec R_i - \vec r) \times (\vec R_{i+1} - \vec r) =
  \vec R_i \times \vec R_{i+1} - \vec r \times (\vec R_{i+1} - \vec R_i)
\end{equation*}

$\vec R_i \times \vec R_{i+1}$ and $\vec R_{i+1} - \vec R_i$ are both fixed
for all time (unique to each vertex), so this approach avoids
computing the vertex-to-point vectors for each vertex.

%******************************************************************************
\subsection{Polygon Enclosure (Unit Sphere)}
Much like the planar case, there are three phases to identify whether a
point lies within a bounding convex polygon:
\begin{itemize}
  \item Confirm that the vertices make a convex polygon.
  \item Bound the polygon with a simple min/max box to streamline
    identification of points that are clearly outside the polygon.
  \item Evaluate whether the point lies within the simple min/max bounds and,
    if it does, evaluate whether it lies within the polygon.
\end{itemize}

\subsubsection{Testing for Convex}

To verify the convex nature, we compare the vector-products between adjoining
line segments as we did for the planar case. However, there are two
important considerations that distinguish the spherical case from the planar
case:
\begin{itemize}
  \item In the planar case, the vertices were all in the same Euclidean plane;
    the vector-product between adjacent edges all had a component along the
    vector perpendicular to that plane (the Cartesian z-axis), which was
    consistent for all polygon edges.
    In the spherical case, the vertices are distributed in three dimensions
    and the polygon formed from the Euclidean straight-line segments
    between vertices is, in general (for more than 3 vertices), a skew
    polygon not a planar polygon.
  \item Because the vertices are restricted the unit spherical surface, the
    polygon can still be described in two non-Euclidean dimensions,
    but the polygon edges
    must be arcs of the great-circles passing through adjacent vertices.
    The radial vector to the surface is analogous to the z-axis from the
    planar case; however, it is important to recognize that, when working
    in Cartesian coordinates, the radial vector varies from vertex-to-vertex.
\end{itemize}


Consider a polygon comprising $N$ vertices, $V_i$, $i \in [1,N]$ on a
unit-sphere with center at point $O$. For each pair of adjacent vertices,
we can define a plane with point $O$; the intersection of this plane with
the unit-sphere surface defines the great-circle line segment between the
two vertices. Each vertex $V_i$ has a corresponding radial vector
$\vec R_i$. The plane defining the polygon line segment between $V_i$ and
$V_{i+1}$ is defined by the perpendicular vector
$\vec S_i = \vec R_i \times \vec R_{i+1}$. The convex nature of the polygon
can be evaluated by taking the vector product of adjacent vectors
$\vec S_{i-1} \times \vec S_{i}$ and comparing the direction against the
radial vector $\vec R_i$ to the common vertex $V_i$.

\begin{equation*}
  \alpha = \left( \vec S_{i-1} \times \vec S_{i} \right) \cdot \vec R_i
\end{equation*}
If this value has the same sign for every vertex, the polygon is convex. For
example, if this value is negative, the two $\vec S$ vectors point ``outward''
and the polygon progresses clockwise as viewed from outside the sphere;
if the value is negative for \textit{every} vertex, then all $\vec S$
vectors point ``outward'' and the polygon is therefore convex.

This can be simplified using vector quadruple and triple product rules:
\begin{align*}
 \alpha & = \left( \left( \vec R_{i-1} \times \vec R_{i} \right) \times
                   \left( \vec R_{i} \times \vec R_{i+1} \right)
            \right)
            \cdot \vec R_i \\
 {} &= \left( \left( \vec R_{i-1} \cdot
                     \left( \vec R_{i} \times \vec R_{i+1} \right)
              \right) \vec R_i -
              \left( \vec R_{i-1} \cdot
                     \left( \vec R_{i} \times \vec R_{i} \right)
              \right) \vec R_{i+1}
       \right) \cdot \vec R_i  \\
 {} &= \left( \vec R_{i} \cdot \left( \vec R_{i+1} \times \vec R_{i-1} \right) -
              [0]
       \right) \vec R_i \cdot \vec R_i  \\
 {} &= \left( \vec R_{i} \cdot
              \left( \vec R_{i+1} \times \vec R_{i-1} \right)
       \right)
\end{align*}

Notes:
\begin{itemize}
  \item The value of $\alpha$ is not truly important, only its sign is
    significant. This is represented by \textit{direction\_sign} in the code.
  \item For simpler evaluation of scalar- and vector-products,
    the coordinates of the vertices are represented in three-dimensional
    Cartesian coordinates. The vertices may be specified in either
    two-dimensional angular coordinates or three-dimensional Cartesian
    coordinates, but for the inner working of the model, any provided
    angular coordinates will be internally converted to Cartesian coordinates.
\end{itemize}


\subsubsection{Bounding the Polygon}

Bounding the polygon using a Cartesian-coordinate representation for
the vertices is more challenging for the spherical geometry than for the
planar geometry because of the curvature of the polygon edges, and because
the bounding box must be a three-dimesnional right-rectangular prism, rather
than a planar rectangle.

Creation of the bounding box is non-trivial and may not be a valuable
endeavor in all situations. For cases where few points will be rejected for
being outside the bounding box (either because few points will be tested
for each definition of vertices, or because the bounding box encloses most of
the sphere) it may be more efficient to simply test
all points against the polygon edges that to bound the polygon. The model
provides a Boolean \textit{apply\_bounding\_box} flag that may be set to
\textit{false} to circumvent the generation and application of the bounding box.

The process for defining the upper and lower limits of the bounding box is
symmetric for all three axes. This process seems more intuitive for the
z-axis (possibly influenced by our visualization of long-range aircraft
trajectories), so the z-axis is used for illustrative purposes. Furthermore,
identification of the lower bound on each axis follows a similar process as
the indentification of the upper bound; only the upper bound is presented
here.

\begin{itemize}
  \item We start by testing whether the polar vector $[0,0,1]$ lies within the
    polygon. If it does, the upper limit of the bounding box is the pole. We
    set the upper limit to $2$ (being an easily-identifiable  value larger
    than $1$) in this case.

    Note -- for purposes of testing whether the pole lies within the polygon,
    we have to switch off the \textit{apply\_bounding\_box} flag to prevent
    the pole from inadvertently being rejected for lying outside of an
    as-yet-undefined bounding box.

  \item If the pole is not enclosed within the polygon, the next step is to
    find the maaximum z-value from the set of vertices. The bounding box
    must include all of these values, so this is a first -- but incomplete
    -- step.

  \item In general, the polygon edges will extend beyond the z-values of the
    vertices, due to the polygon edges being great-circles between the
    vertices.

    Consider, as a simple illustrative example, two mid-latitude
    vertices separated in longitude by 180 degrees. The great-circle passing
    through these points passes through the pole, so the polygon -- and
    therefore the bounding box -- must enclose the pole. Capping the
    bounding box at the highest vertex z-coordinate will result in rejecting
    points for being outside the box that are within the polygon.

    Consequently, we need to evaluate the peak z-value of the polygon edges
    as well as the peak z-values of each vertex. There are some exceptions:
    \begin{itemize}
      \item If all vertices have z-values $< 0$, there are no edges that could
        set a new maximum z-value above that of the value set by the vertices.
      \item Only those edges along the upper part of the polygon can set a
        new maximum z-value. These edges have a $\vec S$ vector with a
        positive z-component if $\alpha <0$ or a negative z-component if
        $\alpha >0$. In either case, $\alpha \vec S_{i,z} < 0$ for suitable
        edges.
      \item It may also be possible, under
        certain polygon construction limitations, to assume that the edge
        that reaches the highest z-value is one of the two edges associated
        with the vertex at the highest z-value, and limit this search to just
        two edges. However, this is not universally true so is not assumed in
        the baseline implementation.
    \end{itemize}
    The algorithm for computing the peak z-value of an edge follows.
\end{itemize}

Consider two adjacent vertices $V_i$ and $V_{i+1}$ and their corresponding
plane-vector $\vec S_i = \vec R_i \times \vec R_{i+1}$, as defined above.
The term \textit{vertex-plane}  is used to identify the plane perpendicular
to vector $\vec S_i$.
Recall, the intersection of this vertex-plane with the unit spherical
surface defines one polygon edge.

The vector passing through the peak z-values is perpendicular to
\begin{itemize}
  \item The vector defined by the intersection of the vertex-plane and the
    x-y plane; that vector is perpendicular to $\vec S_i$ and $\hat k$.
  \item $\vec S_i$, as are all vectors in the vertex-plane.
\end{itemize}
Call this vector $\vec T_i$.

\begin {equation*}
\vec T_i = \left( \vec S_i \times \hat k \right)\times \vec S_i
\end {equation*}

Note that the code uses this general form; it may be further simplified as
follows, but the simplified form is more difficult to generalize to multiple
axes and directions.
\begin {align*}
\vec T_i 
       &= (\vec S_i \cdot \vec S_i) \hat k - (\vec S_i \cdot \hat k) \vec S_i \\
       &= [-S_{i,x} S_{i,z}, -S_{i,y} S_{i,z}, (S_{i,x}^2 + S_{i,y}^2)]
\end {align*}


The polygon edge is some arc of that great circle which may or may not
include the the peak z-value identified by $\vec T_i$.
Note that the peak z-value forms a turning-point; if the arc does not
include a turning point then it is monotonic and bounded by its ends,
which are the vertices. We have already
evaluated the peak z-value from the vertices, so we need only consider the
peak value of the great circle if $\vec T_i$ lies between $\vec V_i$ and
$\vec V_{i+1}$ in the vertex-plane. This can be tested by evaluating
$\left( \vec V_i \times \vec T \right) \cdot \vec S_i$.

Because $\vec T_i$ lies in the vertex-plane, $\vec V_i \times \vec T$ must be
aligned with $\vec S_i$.

Thus $\left( \vec V_i \times \hat T \right) \cdot \vec S_i$ is a measure of the
angle in the vertex plane between $\vec V_i$ and $\vec T_i$. Similarly,
$\left( \vec V_i \times \vec V_{i+1} \right) \cdot \vec S_i =
\lVert\vec S_i \rVert$ is an equivalent measure of the angle between
$\vec V_i$ and
$\vec V_{i+1}$.

Thus:
\begin{itemize}
  \item $\left( \vec V_i \times \vec T \right) \cdot \vec S_i > 0$ implies the
    angle to $\vec T_i$ is in the same direction as that to $\vec V_{i+1}$, and
  \item $\left( \vec V_i \times \vec T \right) \cdot \vec S_i <
    \lVert\vec S_i \rVert \lVert \vec T_i \rVert$ implies the angle to $\vec T_i$
    is smaller than that to $\vec V_{i+1}$
\end{itemize}
With both conditions satisfied, we know the peak z-value from the great circle
lies between the two vertices and we need to consider this value as a potential
maximum z-value for the bounding box..

Normalizing $\vec T_i$ (i.e. to unit magnitude) provides the Cartesian
coordinates of the extreme-latitude points of the great-circle, so the peak
z-value obtained by the great circle is $(\hat T)_z$.


\subsubsection{Determining Whether a Point is Within the Polygon}

As with the planar case, a point $P$ lies within the polygon edges if the
vector-product of adjacent vertex-to-P vectors all point in the same general
direction. For the planar case, this general direction was the Cartesian
z-axis. For the spherical case, it is the radial vector associated with
point $P$.

For each vertex $V_i$ with position vectors $\vec R_i$, we can set a vector
from the vertex to the point of interest,
$\vec r_i = \vec R_P - \vec R_i$. For the point to be within the polygon,
$\forall i, \beta_i  = (\vec r_i \times \vec r_{i+1}) \cdot \vec R_P$
has the same sign.

That scalar-product actually simplifies nicely: as with the planar case,
we don't need to difference the vectors to the vertex and the point
for each vertex because:
\begin{align*}
\beta_i &=
     (\vec r_i \times \vec r_{i+1} ) \cdot \vec R_P) \\
{}&= ((\vec R_P - \vec R_i) \times (\vec R_P - \vec R_{i+1})) \cdot \vec R_P \\
{}&= (-(\vec R_P \times \vec R_{i+1}) - (\vec R_i \times \vec R_P ) +
        (\vec R_i \times \vec R_{i+1} )) \cdot \vec R_P \\
{}&= (\vec R_i \times \vec R_{i+1} )) \cdot \vec R_P
\end{align*}

As with the planar case, we can simplify the identification of whether a
point lies within the polygon by computing and reusing a value that is a
function of the geometry of the polygon, and then including the coordinates of
the point as a simple calculation with this value.

The spherical symmetry produces a result simpler than the planar case. For
this spherical case, we need only the cross product between adjacent
vertices.

Additional notes:
\begin{itemize}
\item If point $P$ lies in any vertex-plane, $\beta = 0$ for that plane.
  Point $P$ lies in that plane if and only if it lies on the polygon edge.
  This is considered within the polygon.
\item The analysis above was performed using Euclidian geometry, but applies
  equally to a spherical geometry. Because the polygon is convex (already
  demonstrated), when dealing with spherical geometry we could evaluate the
  vector-product of the Euclidean vectors representing the
  perpendiculars to the planes defined by the vertex $V_i$ and points $P$ and
  $O$ (instead of the vector products of the vectors $r_i$).
  The vector-product of these vectors simplifies down to the same result.
\end{itemize}

%******************************************************************************
%******************************************************************************
%******************************************************************************
%******************************************************************************
\chapter{User's Guide}

\subsection{Polygon Enclosure (Planar)}
The Polygon Enclosure Method is best implemented as a class member; the
capability is implemented as a class-template, requiring the number of
vertices:
\begin{code}
#include "cml/models/utilities/polygon_enclosure/include/polygon_enclosure.hh"
...
PolygonEnclosure<6> enclosure;
\end{code}

The constructor for the model takes 2 arguments:
\begin{enumerate}
\item A reference to the x-coordinate of the point of interest.
\item A reference to the y-coordinate of the point of interest.
\end{enumerate}

\begin{code}
enclosure ( x, y)
\end{code}

The polygon is configured with the \textit{set\_xy(...)} method.
The argument passed to \textit{set\_xy(...)} is a $N \times 2$ C-style array.
\begin{code}
double enclosure_vertices[6][2] = {...};
...
enclosure.set_xy{ enclosure_vertices);
\end{code}

The model must be initialized prior to use, and after the vertices have been
set. This checks the polygon for being convex and establishes the
rectangular box around the polygon.
\begin{code}
enclosure.initialize();
\end{code}

To test whether the point at coordinates provided at construction lies
within the polygon defined with the coordinates provided by
\textit{set\_xy(...)}, call \textit{in\_polygon}.
\begin{code}
bool within_polygon = enclosure.in_polygon();
\end{code}


%******************************************************************************
\subsection{Polygon Enclosure (Unit Sphere)}
The Polygon Enclosure Method is best implemented as a class member; the
capability is implemented as a class-template, requiring the number of
vertices:
\begin{code}
#include "cml/models/utilities/polygon_enclosure/include/polygon_enclosure_sphere.hh"
...
PolygonEnclosureSphere<6> enclosure;
\end{code}

The polygon is configured with the \textit{set\_vertices(...)} method.
The argument passed to \textit{set\_vertices(...)} can be either:
\begin{itemize}
  \item a $N \times 2$ C-style array representing the $[\lambda, \phi]$ coordinates
    for each vertex.
  \item a $N \times 3$ C-style array representing the $[x,y,z]$ coordinates
    for each vertex.
\end{itemize}

\begin{code}
double enclosure_vertices[6][2] = {...};
...
enclosure.set_vertices{ enclosure_vertices);
\end{code}

The model must be initialized prior to use, and after the vertices have been
set. This checks the polygon for being convex and establishes the
rectangular box around the polygon.
\begin{code}
enclosure.initialize();
\end{code}

To test whether the point at coordinates provided at construction lies
within the polygon defined with the coordinates provided by
\textit{set\_vertices(...)}, call \textit{in\_polygon(...)}.
The argument passed to \textit{in\_polygon(...)} can be either:
\begin{itemize}
  \item a 2-element array representing the $[\lambda, \phi]$ coordinates of
     the point.
  \item a 3-element C-style array representing the $[x,y,z]$ coordinates
    of the point.
\end{itemize}


\begin{code}
bool within_polygon = enclosure.in_polygon( R_point);
\end{code}


%******************************************************************************
%******************************************************************************
%******************************************************************************
%******************************************************************************
\chapter{Inspection, Verification and Validation}
TBD
\end{document}
